%%%%%%%%%%%%%%%%%%%%%%%%%%%%%%%%%%%%%%%%%%%%%%%%%%%%%%%%%%%%%%%%%%%%%%%%%%%%%%%%
% Official RA-L initial-submission style, advisor-facing outline version.
% Source template: RA-L author page -> PaperCept ieeeconf.zip.
%%%%%%%%%%%%%%%%%%%%%%%%%%%%%%%%%%%%%%%%%%%%%%%%%%%%%%%%%%%%%%%%%%%%%%%%%%%%%%%%

\documentclass[letterpaper, 10 pt, conference]{ieeeconf}

\IEEEoverridecommandlockouts
\overrideIEEEmargins

\usepackage{amsmath,amssymb}
\usepackage{graphicx}
\usepackage{booktabs}
\usepackage{xcolor}

\newcommand{\method}{ToolFrame}
\newcommand{\todo}[1]{\textcolor{red}{TODO: #1}}
\newcommand{\claim}[1]{\textbf{Claim:} #1}
\newcommand{\openitem}[1]{\textit{Open:} #1}

\title{\LARGE \bf
\method{}: Let the Tool Define the Frame for Dexterous Manipulation
}

% RA-L initial submissions are double-anonymous.
\author{}

\begin{document}

\maketitle
\thispagestyle{empty}
\pagestyle{empty}

\begin{abstract}
This outline frames \method{} around a bottom-level mechanism: tool use succeeds
only when the functional geometry of a tool is made to coincide with the
desired physical effect in the scene. The difficulty is that this mechanism is
not directly given by language, object category, or a VLM-generated trajectory.
It must be recovered from public evidence, expressed in metric geometry, and
executed by a dexterous policy without relying on hidden evaluator references
or task-specific fallback primitives. \method{} addresses this by separating
tool-use reasoning into three explicit layers: sparse visual evidence
selection, tool-centric geometric compilation, and fixed-goal dexterous
rollout. The key paper message is that the VLM should not be the trajectory
generator; it should expose the visual evidence from which a compiler can bind
tool contact, tool motion, and scene effect into whole-tool $SE(3)$ subgoals.
\todo{Final abstract should add frozen headline results only after the
benchmark is finalized.}
\end{abstract}

\section{INTRODUCTION}

\subsection{Bottom-Level Mechanism}
The paper should begin from the physical mechanism of tool use rather than from
VLM planning. A tool is not useful because it belongs to a semantic category;
it is useful because one part of its geometry transfers contact, force, or
motion to the scene. A brush acts through bristles, an eraser through a flat
patch, a hammer through an impact face, a screwdriver through a tip and axis,
and a marker through a point contact. In dexterous manipulation, the hand may
grasp the same tool in many ways, but the task succeeds only when the whole
tool is placed so that the functional part and functional direction create the
desired scene effect.

\subsection{Core Difficulty}
The fundamental challenge is a three-way binding problem. The system must bind
the tool mechanism, the scene mechanism, and the dexterous execution interface:
\begin{equation}
  \begin{aligned}
  &\text{tool mechanism}\\
  &\quad \longleftrightarrow \text{scene effect}
  \longleftrightarrow \text{whole-tool execution}.
  \end{aligned}
\end{equation}
Existing VLM-centric robot pipelines often blur this binding by asking the VLM
to produce actions, dense waypoints, or trajectories. That is attractive
engineering-wise, but weak scientifically: success can come from hidden target
priors, task-family assumptions, controller behavior, or prompt-specific
shortcuts rather than from image-grounded tool-use reasoning.

\subsection{Central Insight}
The main insight is to treat tool use as a \emph{compiler problem}, not a
trajectory-generation problem. The VLM contributes sparse public references:
which tool part is functional and which scene points or directions matter. The
compiler is responsible for turning these references into metric constraints on
the whole tool. The policy is responsible only for executing the resulting
fixed goals. This division makes the causal chain visible:
\begin{quote}
public evidence $\rightarrow$ functional references $\rightarrow$ whole-tool
$SE(3)$ subgoals $\rightarrow$ dexterous rollout $\rightarrow$ task contract.
\end{quote}

\subsection{Proposed Positioning}
\claim{\method{} is an image-grounded, tool-centric pipeline for dexterous tool
use where the tool defines the manipulation frame and the hand follows
compiled whole-tool subgoals.}

The work should be positioned against two common but weaker framings. It is not
``VLM writes robot trajectories'', because the VLM only selects sparse public
evidence. It is not merely ``affordance prediction'', because the predicted
tool references become geometric constraints that are compiled into executable
subgoals. The novelty is the explicit bridge between functional tool evidence,
RGB-D scene evidence, and fixed-goal dexterous execution.

\subsection{Contributions to Claim}
\begin{enumerate}
  \item \textbf{Mechanism-explicit formulation.} Tool use is formulated as
  aligning a tool-side functional mechanism with a scene-side effect mechanism.
  \item \textbf{Evidence-bounded VLM role.} The VLM selects sparse public tool
  and scene references, but does not generate trajectories, primitives, or
  task-family fallbacks.
  \item \textbf{Tool-centric geometric compiler.} The compiler converts
  contact and motion references into whole-tool $SE(3)$ subgoals for a
  dexterous fixed-goal policy.
  \item \textbf{Task-contract evaluation.} The benchmark evaluates final
  physical effect and attributes failures along the mechanism chain.
\end{enumerate}

\section{RELATED WORK POSITIONING}

\subsection{VLMs for Robot Planning}
The relevant prior work shows that VLMs are useful semantic and visual
interfaces. The gap is that free-form VLM plans do not expose the geometric
mechanism needed for contact-rich tool use. Our distinction should be stated
crisply: \method{} uses VLMs for evidence selection, while metric geometry and
execution are handled by explicit modules.

\subsection{Tool Affordances and Functional Parts}
Affordance and part-grounding work often identifies where a tool can be used.
The missing step is converting this knowledge into robot-executable constraints.
\method{} makes the functional part operational by representing it through a
contact point and motion axis that are consumed by a compiler.

\subsection{Constraint-Based Manipulation}
Constraint-based manipulation provides the right language for geometric
reasoning, but usually assumes task frames or contacts are already specified.
\method{} contributes the image-grounded front end and tool-centric reference
interface that produce those constraints from public evidence.

\subsection{Benchmarks for Tool Use}
The benchmark contribution should be framed around attribution, not only around
success rate. The paper should show not just whether a rollout succeeded, but
where the mechanism chain failed when it did not.

\section{PROBLEM STATEMENT}

\subsection{Mechanistic Variables}
At the outline level, the technical interface can stay compact:
\begin{equation}
  \mathcal{R}_{tool} =
  \{p^{contact}_{tool}, a^{motion}_{tool}\},
\end{equation}
\begin{equation}
  \mathcal{R}_{scene} =
  \{p^{start}_{scene}, p^{path}_{scene}, p^{end}_{scene},
  a^{direction}_{scene}\}.
\end{equation}
The paper's technical problem is to infer these references from public
evidence and compile them into whole-tool poses
$T_{world,tool}\in SE(3)$ that are executable by the fixed-goal policy.

\subsection{Evidence Boundary}
This boundary is part of the paper's scientific value. The method may use
public RGB-D scene evidence, public tool-part evidence, tool meshes, task
language, and robot state needed for execution. It must not use hidden
evaluator targets, oracle target axes, oracle trajectories, primitive
fallbacks, repair loops, task-family fallbacks, or seed-specific shortcuts.

\subsection{Why This Boundary Matters}
The boundary prevents the benchmark from becoming a test of hidden priors. If a
method can snap to simulator targets or fall back to hand-coded task
primitives, then the central claim of image-grounded tool use is no longer
testable. \method{} should be sold as a stricter, more auditable route: it uses
less privileged information and exposes the intermediate mechanism.

\section{METHOD OUTLINE}

\subsection{Overview Figure}
Fig.~\ref{fig:overview} should be the main selling figure. It should not read
like a software pipeline. It should read like a mechanism decomposition:
tool-side mechanism, scene-side mechanism, compiler bridge, and dexterous
rollout.

\begin{figure*}[t]
\centering
\fbox{\parbox{0.92\textwidth}{\centering
Advisor-facing Fig. 1 placeholder: (A) recover tool mechanism;
(B) ground scene effect references; (C) compile whole-tool $SE(3)$ subgoals;
(D) execute and evaluate with task contracts.}}
\caption{Conceptual overview. \method{} exposes the hidden mechanism of tool
use: public tool evidence defines contact and motion references; public RGB-D
scene evidence defines effect references; a geometry compiler binds both into
whole-tool subgoals; a fixed-goal dexterous policy executes them.}
\label{fig:overview}
\end{figure*}

\subsection{A. Tool-Centric Prior Extraction}
\claim{Recover the functional mechanism of the tool.}

This section should explain why object category and visible geometry are
insufficient. The method extracts a compact tool-side representation: a
functional part, a contact point, and a motion axis. Avoid over-specifying the
current implementation details unless they are frozen. The important point is
that the output is not a semantic label but a geometric mechanism.

\subsection{B. Scene Key-Points Grounding}
\claim{Recover the sparse scene mechanism of the desired effect.}

This section should explain that the task does not require the VLM to draw a
robot trajectory. It requires sparse references that describe where the
physical effect should occur. RGB-D lifting turns selected image refs into
metric scene refs. Keep uncertain details such as exact prompting, sampling
policy, or per-task thresholds out of the main outline until they are stable.

\subsection{C. Geometry-Constrained Compiler}
\claim{Bind the tool mechanism to the scene mechanism.}

This is the most important method section. The compiler should be presented as
the bridge that turns public evidence into executable whole-tool poses. At the
outline level, state the constraint families rather than implementation
particulars: contact alignment, motion-axis alignment, approach/clearance, and
stage-wise whole-tool subgoals. The compiler's output is fixed $SE(3)$ goals,
not low-level robot actions.

\subsection{D. Fixed-Goal Dexterous Rollout and Evaluation}
\claim{Let the hand execute; let task contracts judge the physical effect.}

The policy section should be deliberately modest. The method does not claim to
learn a new universal dexterous policy. It uses a fixed-goal SimToolReal policy
as the execution layer for compiled subgoals. The evaluation layer then
attributes failures to grounding, compilation, execution, or final task effect.

\section{BENCHMARK AND EVALUATION PLAN}

\subsection{What the Benchmark Should Demonstrate}
The benchmark should demonstrate that the formulation is tool-centric rather
than task-name-specific. The strongest version would cover multiple tool
geometries and at least two scene families, while clearly separating headline
statistics from candidate or diagnostic tasks.

\subsection{Task Tiers}
\begin{table}[t]
\caption{Advisor-facing task tiers. Exact task membership should be frozen only
after contracts and launcher summaries are finalized.}
\label{tab:tiers}
\begin{center}
\begin{tabular}{ll}
\toprule
Tier & Role in the paper \\
\midrule
Headline & Main statistics from frozen contracts \\
Candidate & Promising tasks pending final review \\
Diagnostic & Stress tests for mechanism failures \\
Draft/Smoke & Development evidence, not final claims \\
\bottomrule
\end{tabular}
\end{center}
\end{table}

\subsection{Metrics}
Use metrics to explain the mechanism chain, not only to report success:
functional grounding, scene grounding, subgoal validity, policy execution,
final task success, and first-failure attribution. Do not insert numbers until
the frozen launcher summaries are available.

\subsection{Expected Main Table}
The final paper should contain one compact table for headline tasks:
\begin{quote}
Task, Tool, Scene, $N$, grounding validity, subgoal validity, execution
completion, final success, and dominant failure mode.
\end{quote}
The exact columns can be adjusted after the contracts stabilize.

\section{PAPER PACKAGING}

\subsection{One-Sentence Pitch}
\begin{quote}
\method{} makes dexterous tool use auditable by turning public tool and scene
evidence into tool-centric whole-tool subgoals, instead of asking a VLM to
invent robot trajectories.
\end{quote}

\subsection{Main Figure Message}
The figure should visually prove the paper's story:
\begin{quote}
Tool geometry defines the functional frame; RGB-D scene evidence defines the
effect frame; the compiler binds them; the hand executes.
\end{quote}

\subsection{Claims to Avoid}
The outline should not overclaim general tool use, general dexterity, online
adaptation, real-robot transfer, or open-ended VLM planning unless experiments
support them. It should not report historical migration snapshots as final
success statistics.

\section{OPEN DECISIONS FOR ADVISOR DISCUSSION}

\begin{itemize}
  \item Which title is strongest: the current slogan title or a more
  conservative tool-centric visual grounding title?
  \item Which tasks are safe as headline statistics after contract freezing?
  \item How many tool families and scene families are necessary for the RA-L
  claim to feel broad enough?
  \item Which baseline is most fair: direct VLM waypointing, no tool axis, no
  contact point, or primitive-style task planning?
  \item Should the paper emphasize ``dexterous manipulation'' in the title, or
  reserve dexterity for the execution section to avoid overclaiming?
\end{itemize}

\section{LIMITATIONS}

The paper should state limitations as consequences of the mechanism boundary.
If the functional part is not visible or reviewed, the tool mechanism may be
wrong. If depth lifting is noisy, the scene mechanism may be wrong. If the
compiler produces a feasible tool pose but the policy cannot maintain contact,
the task may still fail. These limitations are acceptable because the paper
does not hide them behind repair loops or privileged simulator references.

\section{CONCLUSION}

The paper's final message should be simple and defensible: tool use is a
mechanism-binding problem. \method{} exposes this mechanism through public
visual evidence, binds tool and scene references with geometry, and executes
the result through a fixed-goal dexterous policy. This is the strongest
packaging of the work: the tool defines the frame, the compiler makes it
executable, and the hand follows.

\bibliographystyle{IEEEtran}
\bibliography{references}

\end{document}
